% ===========================================================================
%  A blank LOOM fill-in notebook.  Copy this folder and start writing.
%  Compile:  latexmk -xelatex main      (or: xelatex main, twice)
%  Options:  [cjk] 中文 · [letter] US-letter paper
%  Everything below is a live cheat-sheet of what the class gives you.
% ===========================================================================
\documentclass[cjk]{loom}

% math shorthands you want (these are per-document):
\newcommand{\R}{\mathbb{R}}
\newcommand{\E}{\mathbb{E}}

\runningthread{your topic}          % quiet footer label

\begin{document}

\loomcover
  {Title of the Notebook}
  {a one-line subtitle}
  {Your Name}
  {\today}

\section{First course}
\warmth{0}\quad\whisper{One-line orientation for this section.}

\begin{strand}
The \keyword{big idea}, in your own words, before the machinery. This is the
\emph{passive} layer: written to be read.
\end{strand}

\block{The paradigm}
\trigger{When you'd reach for this.}
A displayed master fact:
\[ \text{(your key equation here)} \]

\block{Cheat-sheet}
\noindent\begin{tabularx}{\linewidth}{@{}l l L@{}}
\toprule
{\color{indigo}Tool} & {\color{indigo}Setting} & \multicolumn{1}{l}{\color{indigo}When to use}\\
\midrule
A & a setting & a wrapping description that flows to the next line if it is long\\
B & another   & another note\\
\bottomrule
\end{tabularx}

\block{Results (read the statement, fill the proof)}

\begin{definition}[a name]\label{def:thing}
A madder knot. State the object, but leave the key clause blank: closed under
\fillin[3cm] and \fillin[2cm].
\end{definition}

\begin{theorem}[attribution]\label{thm:main}
An indigo knot. State the result in full.
\end{theorem}
\begin{proof}
Skeleton: (i) the first move; (ii) the second. \TODO{the step that makes it work}
\end{proof}

\begin{example}
A weld knot for a worked instance.
\end{example}

\begin{yourturn}
Restage the example as a computation you do. Show that \dots
\workspace[3]   % three ruled lines to write on (delete if you type your answer)
\end{yourturn}

\begin{remark}
An unboxed aside.
\end{remark}

\recall{A quick active-recall question, parked in the margin.}%
Inline prose continues here. \warmth{2} marks how well you grok this block (0--5).

\loose{An open thread / exercise to pull next time.}%
% recurring threads:  \warp{key} declares one, \pick{key} cross-refs it later.

\end{document}
